\documentclass[11pt]{article}

\usepackage{sectsty}
\usepackage{graphicx}

% Margins
\topmargin=-0.40in
\evensidemargin=0in
\oddsidemargin=0in
\textwidth=6.5in
\textheight=9.0in
\headsep=0.25in

\title{ Focus Lab}
\author{ Pietro Domingues }
\date{\today}

\begin{document}
\maketitle	
\pagebreak



% Optional TOC
% \tableofcontents
% \pagebreak

%--Paper--

\section{Introdução}

Focuslab nasceu com o propósito de melhorar o processo cognitivo e o conforto da mente.\\

Não queremos nos categorizar como uma solução de "produtividade" num sentido de que "precisamos ser produtivos". Somos uma solução para ajudar a sua mente a focar \textit{no que importa.}\\

O que importa pode ser desde seu trabalho e sua produção, até mesmo a sua "missão", sua vocação ou seu "sentido espiritual". E o que importa varia com o tempo e situação, pois somos multi-facetados.

\section{Princípios}


\section{Referências e estado da arte}

\subsection{Applicativos existentes}

\begin{itemize}
\item Opal: Ótimo para controlar atenção e bloquear distrações. Não muito intuitivo para substituir com tarefas interessantes.
\item Toggl Track: Ótimo para trackear tempo e projetos.
\item Todoist: Ótimo para resolver tarefas e ter metas.
\end{itemize}

\subsection{Técnicas existentes}

\begin{itemize}
\item Get Things Done (GTD): Arcabouço organizacional para ser all-in-one. Funciona por um tempo mas pode ficar muito repetitivo.
\item Atomic Habits: O core da iniciativa é muito boa e se inspira em práticas yogicas em algum sentido
\item Yoga: O foco e atenção plena é a principal coisa que interessa para podermos portar para um aplicativo voltado a foco e atenção plena.
\item TCC: Considerar condicionamentos é uma parte restrita mais valiosa em determinados tipos de tarefas.
\end{itemize}

\subsection{Referências e corpo teórico}

\begin{itemize}
\item Psicanálise: 
\item Psicologia analítica:
\item Neurociência:
\item Yoga:
\item Teoria da informação:
\item 
\end{itemize}

\section{Topologia do projeto}

O projeto pode ser separado em pastas, cada uma com um propósito.
\begin{itemize}
    \item Monorepo focuslab: O projeto principal, monorepo, com todas a base de informações do focuslab.
    \item Focus: O aplicativo principal, estabilizado e com ideias maduras, prontas para ser utilizadas com pessoas reais.
    \item Lab: O laboratório de ideias, conceitos, protótipos para serem experimentados livreamente sem ter restrição de funcionar com estabilidade.
    \item Documents: Os documentos e referências que guiam a pesquisa como um todo.
\end{itemize}

\section{Desenvolvimento}

Stack de desenvolvimento:

\section{Diário de bordo}

\subsection{17/12 - Quarta-feira}

Preciso fazer uma compilação de ideias novamente e continuar do ponto onde parei.\\

Nutrir os próximos passos com referências do dia a dia.\\

Reservar 1 hora por dia para fazer qualquer coisa relacionada ao focuslab.

\subsection{18/12 - Quinta-feira}

Bom dia! Hoje acordei animado e estive pensando em algumas coisas ontem à noite antes de dormir...\\

Uma delas é que o aplicativo é seu e você tem liberdade para voltar para o que você quiser! Contanto que haja foco e disciplina na execução, as ideias podem se moldar à vontade para as suas necessidades imediatas.\\

Limitar tempo e priorizar. Constantemente.\\

Refleti ao final do dia: A natureza tem um tempo próprio. Se observamos as estruturas que chamamos de harmônicas e sincrônicas se devem ao fato de haverem múltiplas interações em um campo de tensões contínuo. E essas estruturas se acomodam não com pressa, mas de acordo com os tempos próprios e assim confluem como um mecanismo autônomo para algo integral.\\

O propósito unidirecionado que vem de dentro é o motor mais potente dos seres humanos. Mais potente, duradouro e estruturante do que qualquer outro artíficio externo, qualquer pressão ou alienação de fora.\\


\subsection{19/12 - Sexta-feira}

Hoje foi um dia daqueles... Mas estou muito pensativo em relação à finalidade deste aplicativo. Eu sei que preciso realizar algo.\\

O primeiro pensamento que vem é justamente um MVP auto-referente. Algo que eu consiga utilizar agora para ajudar a desenvolver esse aplicativo da mesma maneira livre, mas organizada.

\subsection{20/12 - Sábado}

-

\subsection{21/12 - Domingo}

Aqui é o primeiríssimo protótipo do que será o futuro do focuslab. \\

Pensei qual seria o MVP para manter o próprio focuslab vivo, e o próprio diário de bordo do focuslab, de forma muito simples, como texto é o primeiro MVP.\\

Não preciso como requisito um aplicativo elaborado onde ainda tenho que gastar energia e tempo (que não tenho às vezes), para validar algumas ideias. Mero texto, disciplina e atenção próprias minhas é que serão a fonte para o MVP estar vivo.\\

O futuro do aplicativo será uma emulação da própria dinâmica deste ser vivo que sou eu com um corpo de texto que por si só é morto. O aplicativo será uma evolução disso, trazendo mais elementos "vivos" para possivelmente qualquer pessoa mesmo que não exiba a princípio uma disciplina ou foco próprio bem elaborados.\\

Mas eu já estou com o foco determinado por mim mesmo, e posso usar de quaisquer ferramentas para manter esse foco vivo. \\

Então a partir de hoje eu declaro aqui que vou utilizar este meio como ferramenta, como MVP para o futuro focuslab.\\

Por enquanto é um pouco rígido e centralizado, mas justamente esse poder que criará um cultivo como cultivamos o templo de nossas relações íntimas com a natureza, Deus e com todas as entidades.\\

Deverá ser guardado de forma sagrada para irradiar cada vez mais para outros lugares, sendo o aplicativo a forma de instrumento e ferramenta que poderá ser utilizada de forma democrática e poderá deter, mesmo que de forma pouca, o avalanche de informações e entropia gerada por uma sociedade sem finalidade.\\

Depois de tudo isso falado então a forma de utilizar este diário vivo será obviamente a partir da própria escrita de um diário, então o compromisso está sendo escrever qualquer coisa que me vem a cabeça a respeito da minha rotina.\\

Então o próximo passo será estabelecer firmezas e compromissos com esse diário, pois caso ele fosse vivo, ele seria capaz de me cobrar e de afetar minha atenção. Pois então esse será meu próximo passo.\\

Sendo assim, eu me comprometo essa semana a fazer:


\begin{itemize}
    \item Um exercício qualquer que seja: caminhada ou trio de exercícios (abdomen, agachamento e flexão).
    \item Um estudo todo dia. Pode ser apenas uma frase, mas preciso pegar um livro da minha biblioteca (é para isso que ela serve). Qualquer ideia vai nutrir de volta meu diário vivo.
    \item Resolver todas as pendências que eu me comprometi. Ser muito claro em dividir pendências do dia e pendências de outros dias. Não se preocupar em pendências restantes pois elas terão seu próprio tempo numa configuração em que você estará mais descansado.
\end{itemize}

\pagebreak
\section{Section 2}
Lorem Ipsum \\

%--/Paper--

\end{document}